% Beeldcredits voor de foto's en AI-illustraties van Boek~5 (Universiteit, derde jaar).
% Machineleesbare licentiegegevens staan in images/book5/CREDITS.md; houd de
% twee bestanden synchroon. De auteur van Boek~5 beheert dit bestand: voeg per
% foto één \item toe (binnen de itemize, per hoofdstuk) zodra foto's binnenkomen.
\chapter*{Bronvermelding van afbeeldingen}

De tekeningen, kaarten en grafieken in dit boek zijn origineel; kaarten en grafieken
die op gepubliceerde gegevens zijn gebaseerd, noemen hun bron in het bijschrift of in
deze lijst. Foto's worden gebruikt onder hun respectievelijke vrije licenties, met
dank aan hun auteurs; volledige bronlinks en licentie-URL's voor elke foto staan in
\texttt{images/book5/CREDITS.md} in de repository van het boek.

\bigskip
Naast de foto's en de originele tekeningen zijn sommige figuren AI-gegenereerde
illustraties, gemaakt met OpenAI-beeldgeneratie op basis van prompts die door de
redacteuren van het boek zijn geschreven, en vóór opname beoordeeld op
wetenschappelijke juistheid. Ze tonen uitsluitend algemene scènes, nooit een echte
plaats, een echt persoon of een echt exemplaar. De volledige prompt voor elke
gegenereerde afbeelding is vastgelegd in \texttt{images/book5/ai/PROMPTS.md} in de
repository.

\bigskip
\noindent Foto's en gegevens, per hoofdstuk:
\begin{itemize}\small
  \item Hfst.~1 (seismologie en aardbeeldvorming): PREM, A.~M.~Dziewonski en D.~L.~Anderson (1981), zoals getabelleerd in ObsPy; ak135, B.~L.~N.~Kennett, E.~R.~Engdahl en R.~Buland (1995).
  \item Hfst.~3 (zwaartekracht en de figuur van de Aarde): de tweelingsatellieten GRACE Follow-On, kunstenaarsimpressie, NASA/JPL-Caltech, publiek domein (Wikimedia Commons). Massabalans van ijskappen van het IMBIE Team (2021), UK Polar Data Centre, Open Government Licence v3.0.
  \item Hfst.~4 (geomagnetisme en paleomagnetisme): International Geomagnetic Reference Field, veertiende generatie, IAGA Division~V Working Group V-MOD, coëfficiënten verspreid door NOAA NCEI.
  \item Hfst.~5 (thermische modellen van de lithosfeer): ouderdom van de oceaanbodem, M.~Seton en anderen (2020), EarthByte, CC BY 4.0; ETOPO1 wereldwijd reliëf, NOAA NCEI, publiek domein.
  \item Hfst.~6 (manteldynamiek en drijvende krachten van plaatbewegingen): de opgeheven stranden van Ballyhillin, Malin Head, A.~F.~Borchert, CC BY-SA 4.0 (Wikimedia Commons); PREM, A.~M.~Dziewonski en D.~L.~Anderson (1981).
  \item Hfst.~7 (deformatiemechanismen en afschuifzones): oogmylonniet, Mikenorton, CC BY-SA 4.0 (Wikimedia Commons).
  \item Hfst.~9 (extensionele en horizontale-verschuivingstektoniek): bekken-en-bergketentopografie ten zuiden van Death Valley, Ken Lund, CC BY-SA 2.0 (Wikimedia Commons).
  \item Hfst.~14 (stabiele-isotopengeochemie): LR04 benthische $\delta^{18}$O-stapel, L.~E.~Lisiecki en M.~E.~Raymo (2005).
  \item Hfst.~15 (water--gesteenteinteractie): travertijnterrassen van Palette Spring, Mammoth Hot Springs, Jim Peaco, National Park Service, publiek domein (Wikimedia Commons). Evenwichtsconstanten uit de PHREEQC-database phreeqc.dat, D.~L.~Parkhurst en C.~A.~J.~Appelo, U.S. Geological Survey, publiek domein.
  \item Hfst.~16 (sequentiestratigrafie): de Book Cliffs boven Helper, Utah, David Mackie, publiek domein (Wikimedia Commons); LR04 benthische $\delta^{18}$O-stapel, L.~E.~Lisiecki en M.~E.~Raymo (2005).
  \item Hfst.~17 (faciesmodellen): hummocky-kruisgelaagdheid in de Arisaig Formation, Michael C.~Rygel, CC BY-SA 3.0 (Wikimedia Commons), bijgesneden.
  \item Hfst.~18 (kwantitatieve geomorfologie): de Niagara-kloof en de Horseshoe Falls, Ken Lund, CC BY-SA 2.0 (Wikimedia Commons), bijgesneden.
  \item Hfst.~19 (geofysische vloeistofdynamica): zonale gemiddelde winden van de NCEP--NCAR Reanalysis~1, gegevens verstrekt door de NOAA PSL, Boulder, Colorado, VS; World Ocean Atlas 2023, NOAA NCEI.
  \item Hfst.~20 (atmosferische dynamica): een lagedrukgebied ten westen van IJsland, NASA/GSFC, MODIS Rapid Response Team, Jacques Descloitres, publiek domein (Wikimedia Commons); zonale gemiddelde winden van de NCEP--NCAR Reanalysis~1, gegevens verstrekt door de NOAA PSL, Boulder, Colorado, VS.
  \item Hfst.~21 (oceaandynamica): een Argo-float uitgezet vanaf een onderzoeksschip, NOAA Global Ocean Monitoring and Observing Program, publiek domein (Wikimedia Commons).
  \item Hfst.~22 (klimaatmodellen): supercomputers van het NASA Center for Climate Simulation, Goddard Space Flight Center, 2010, publiek domein (Wikimedia Commons); oppervlaktelucht\-temperatuur en straling aan de bovenkant van de atmosfeer van de NCEP--NCAR Reanalysis~1, gegevens verstrekt door de NOAA PSL, Boulder, Colorado, VS.
  \item Hfst.~24 (glaciologie): een uitlaatgletsjer in zuidwest-Groenland, NASA / Jim Yungel, Operation IceBridge, 2014, publiek domein (Wikimedia Commons); ijsoppervlakhoogten van Groenland uit ETOPO1 Ice Surface, NOAA NCEI; omvang van het Arctische zee-ijs uit de NSIDC Sea Ice Index, versie~4 (Fetterer en anderen, 2025, doi:10.7265/a98x-0f50).
  \item Hfst.~25 (remote sensing en GIS): het interferogram van 1995--2001 van het vulkanisch centrum Three Sisters, Oregon, U.S. Geological Survey (C.~Wicks en anderen, 2002), publiek domein (Wikimedia Commons); laboratoriumspectra uit de ECOSTRESS Spectral Library, versie~1.0 (S.~K.~Meerdink, S.~J.~Hook, D.~A.~Roberts en E.~A.~Abbott, 2019; A.~M.~Baldridge en anderen, 2009), overgenomen met dank aan het Jet Propulsion Laboratory, California Institute of Technology, Pasadena, Californië.
  \item Hfst.~26 (planetaire binnenkanten en oppervlakken): Barringer Meteor Crater, Arizona, USGS/D.~Roddy, publiek domein (Wikimedia Commons); Olympus Mons, Viking-baanmozaïek, NASA/JPL/USGS, publiek domein (Wikimedia Commons). Krater-productie- en chronologiefuncties zoals gegeven in de Craterstats-software van G.~G.~Michael, naar G.~Neukum, B.~A.~Ivanov en W.~K.~Hartmann (2001).
  \item Hfst.~27 (planetaire atmosferen en bewoonbaarheid): een netwerk van uitgedroogde dalen op Mars, Mars Express HRSC, ESA/DLR/FU Berlin, CC BY-SA 3.0 IGO (Wikimedia Commons); Ligeia Mare, Titan, Cassini-radarmozaïek, NASA/JPL/USGS, publiek domein (Wikimedia Commons). Fits van de bewoonbare zone van R.~K.~Kopparapu en anderen (2013); dampdruk van methaan uit de NIST Chemistry WebBook.
  \item Hfst.~28 (exoplaneten): HIP~65426~b, afgebeeld door de James Webb Space Telescope, NASA/ESA/CSA, A~Carter (UCSC), het ERS~1386-team en A.~Pagan (STScI), CC BY~4.0 (Wikimedia Commons). Bevestigde planeten uit de NASA Exoplanet Archive, tabel Planetary Systems Composite Parameters, geraadpleegd op 23 oktober 2026: dit onderzoek maakt gebruik van de NASA Exoplanet Archive, die wordt beheerd door het California Institute of Technology, in opdracht van de National Aeronautics and Space Administration in het kader van het Exoplanet Exploration Program. Modelkrommen van massa--straal van L.~Zeng en anderen (2016, 2019).
  \item Hfst.~29 (de Melkweg en sterrenstelsels): het panorama van de Melkweg van het GigaGalaxy Zoom-project, ESO/S.~Brunier, CC BY~4.0 (Wikimedia Commons); het spiraalstelsel M101, European Space Agency en NASA (Hubble, heic0602; Hubble-gegevens K.~D.~Kuntz, F.~Bresolin, J.~Trauger, J.~Mould en Y.-H.~Chu; beeldverwerking Davide De~Martin), CC BY~4.0 (Wikimedia Commons). Cirkelsnelheden van de Melkweg van A.-C.~Eilers, D.~W.~Hogg, H.-W.~Rix en M.~K.~Ness (2019); massamodel van E.~Pouliasis, P.~Di~Matteo en M.~Haywood (2017); klassieke Cepheïden van de Grote Magelhaense Wolk uit OGLE, I.~Soszy\'nski en anderen (2015), Acta Astronomica 65, 297, verspreid door het CDS, Straatsburg (VizieR J/AcA/65/297).
  \item Hfst.~30 (waar de elementen vandaan komen): de Krabnevel, NASA, ESA, J.~Hester en A.~Loll (Arizona State University), publiek domein (Wikimedia Commons); de kilonova van GW170817 in NGC~4993, NASA en ESA, dankzij A.~J.~Levan, N.~R.~Tanvir, A.~Fruchter en O.~Fox (Hubble, heic1717a), CC BY~4.0 (Wikimedia Commons). Zonne-abundanties van M.~Asplund, A.~M.~Amarsi en N.~Grevesse (2021); neutroneneinvangstdoorsneden uit KADoNiS v0.3 (I.~Dillmann, R.~Plag, F.~K\"appeler en T.~Rauscher); halveringstijden uit de IAEA Live Chart of Nuclides (ENSDF); metalliciteiten van de Geneva--Copenhagen-survey zoals opnieuw geanalyseerd door L.~Casagrande en anderen (2011), verspreid door het CDS, Straatsburg (VizieR J/A+A/530/A138).
\end{itemize}
